% Generated by scripts/assemble.py; corrections belong in source/corrections.json.
\examhead{70}{大连理工大学 2025 年数学分析试卷}{70}
% p071-b05; source PDF page 71
\par\Needspace{5\baselineskip}
\blockanchor{p071-b05}
\majorquestion{一}简答题. 每题 6 分, 共 60 分.\par
% p071-b06; source PDF page 71
\blockanchor{p071-b06}
\question{1}将 $f(x, y) = \frac{\cos x}{\cos y}$ 在 (0, 0) 展开为麦克劳林泰勒多项式到二阶.\par
% p071-b07; source PDF page 71
\blockanchor{p071-b07}
\question{2}设 $\varphi(x) = \int_{\sin x}^{\cos x} e^{t^2 + xt} \, \mathrm{d}t$, 求 $\varphi'(0)$.\par
% p071-b08; source PDF page 71
\blockanchor{p071-b08}
\question{3}证明 $\zeta(x) = \sum_{n=1}^{\infty} \frac{1}{n^x}$ 在 $(1, +\infty)$ 上连续可微.\par
% p071-b09; source PDF page 71
\blockanchor{p071-b09}
\question{4}证明 $\sum_{n=1}^{\infty} (-1)^{n-1} (\sqrt[n]{n} - 1)$ 条件收敛.\par
% p071-b10; source PDF page 71
\blockanchor{p071-b10}
\question{5}设  $\left|a_{n+p}-a_{n}\right|\leq\frac{p}{n^{2}}, n,p=1,2,\cdots,\left\{a_{n}\right\}$  收敛吗?\par
% p071-b11; source PDF page 71
\blockanchor{p071-b11}
\question{6}当 $x \in (0, \pi)$ 时, 证明: $4(1 - \cos x) < x(x + \sin x)$.\par
% p071-b12; source PDF page 71
\blockanchor{p071-b12}
\question{7}$\lim_{n\to \infty}a_n = A$ 的充要条件是对任意的 $\varepsilon \geq 0$ ，存在正整数 $N$ ，当 $n > N$ 时，有 $|a_{n} - A|\leq \varepsilon$ 是否正确？\par
% p071-b13; source PDF page 71
\blockanchor{p071-b13}
\question{8}$f(x)$ 在 $[a, b]$ 上可微, $f(a) = f(b) = 0$. 证明: 存在 $\xi \in (a, b)$, 使得 $f'(\xi) = 2025f(\xi)$.\par
% p071-b14; source PDF page 71
\blockanchor{p071-b14}
\question{9}$f(x)$ 在 $\mathbb{R}$ 上连续, 以 $T > 0$ 为周期, $\int_0^T f(x)\mathrm{d}x = 0$, $g(x)$ 在 $[0, +\infty)$ 上单调, $\lim_{x\to +\infty}g(x) = 0$. 证明: 广义积分 $\int_0^{+\infty}f(x)g(x)\mathrm{d}x$ 收敛.\par
% p071-b15; source PDF page 71
\blockanchor{p071-b15}
\question{10}证明: $x \sin x$ 在 $\mathbb{R}$ 上不一致连续.\par
% p071-b16; source PDF page 71
\par\Needspace{5\baselineskip}
\blockanchor{p071-b16}
\majorquestion{二}计算题. 每题 10 分, 共 30 分.\par
% p071-b17; source PDF page 71
\par\Needspace{6\baselineskip}
\blockanchor{p071-b17}
\question{1}$\Gamma: y = \sqrt{2x - x^2}$, 从 $(0,0)$ 到 $(2,0)$. 求第二型曲线积分\par
\nopagebreak[4]
% p071-b18; source PDF page 71
\blockanchor{p071-b18}
\displayformula{\int_ {\Gamma} \left(e ^ {x} \sin y - 4 y\right) \mathrm{d} x + \left(e ^ {x} \cos y + 4 x\right) \mathrm{d} y.}
% p071-b19; source PDF page 71
\blockanchor{p071-b19}
\question{2}$V = \{(x,y,z)\mid \sqrt{x^2 + y^2}\leq z\leq 3\}$ ，求 $\iiint_{V}\sqrt{x^2 + y^2}\sin (z^2)\mathrm{d}x\mathrm{d}y\mathrm{d}z.$\par
% p071-b20; source PDF page 71
\blockanchor{p071-b20}
\question{3}求 $\int_0^{+\infty} \frac{e^{-x} - e^{-3x}}{x} \sin x \, \mathrm{d}x$.\par
% p071-b21; source PDF page 71
\par\Needspace{5\baselineskip}
\blockanchor{p071-b21}
\majorquestion{三}证明题. 每题12分, 共60分\par
% p071-b22; source PDF page 71
\blockanchor{p071-b22}
\question{1}设 $f(x)$ 在 $\mathbb{R}$ 上一致连续且有界, 证明: $f_{n}(x) = \frac{n}{1 - e^{-n}} \int_{0}^{1} f(x + t) e^{-nt} \mathrm{d}t$ 一致收敛于 $f(x)$.\par
% p071-b23; source PDF page 71
\blockanchor{p071-b23}
\question{2}$\{a_{n}\}$ 单调递减且 $\lim_{n\to \infty}a_n = 0$ 记 $f(x) = \sum_{n = 0}^{\infty}a_nx^n$ ， $x\in (-1,1)$ ，证明： $\lim_{x\to 1^{-}}(1 - x)f(x) = 0.$\par
% p071-b24; source PDF page 71
\blockanchor{p071-b24}
\question{3}证明在 $(0,0)$ 的邻域内方程 $x - \cos (x^2 y) + e^{x + y^2} = 0$ 决定连续可微函数 $x = x(y)$, 并证明在 $y = 0$ 处取得极大值.\par
% p072-b00; source PDF page 72
\blockanchor{p072-b00}
\question{4}设 $f(x)$ 在 $[0, +\infty)$ 上连续非负，且对任意的 $x, y \geq 0$，有 $f(x + y) \leq f(x) + f(y)$. 证明：$\lim_{x \to +\infty} \frac{f(x)}{x}$ 存在且有限.\par
% p072-b01; source PDF page 72
\blockanchor{p072-b01}
\question{5}设 $f(x)$ 在 $[0, +\infty)$ 上单调递增, 记 $F(x) = \int_0^x f(t) \, \mathrm{d}t (x \geq 0)$. 证明: $F(x)$ 在 $x = 0$ 处存在右导数, 并求 $F_+(0)$.\par
